\documentclass{article}
\usepackage[T1]{fontenc}
\usepackage{lmodern}
\usepackage{graphicx}
\usepackage{amsmath,amssymb}
\usepackage[a4paper,margin=2cm]{geometry}
\usepackage[absolute,overlay]{textpos}
\setlength{\TPHorizModule}{1mm}
\setlength{\TPVertModule}{1mm}
\usepackage[indonesian]{babel}
\usepackage{hyperref}
\setlength{\emergencystretch}{2em}
% unit: Halaman 5

\begin{document}

% === Halaman 5 (python/native) ===

• Pembangkitan kunci putaran sangat sederhana. 

• Kunci eksternal yang panjangnya 256 bit dibagi ke dalam delapan bagian yang 

masing-masing panjangnya 32 bit. 

• Delapan bagian ini yang dinamakan K1, K2, …, K8. 

• Contoh: K = ‘abcdefghijklmnopqrstuvwxyz123456’       (32 x 8 bit = 256 bit)



        maka K1 = ‘abcd’ 


K5 = ‘qrst’



K2 = ‘efgh’ 


K6 = ‘uvwx’



K3 = ‘ijkl’ 


K7 = ‘yz12’



K4 = ‘mnop’ 


K8 = ‘3456’

5

\end{document}
